% Esqueleto didático adaptado do pré-projeto fornecido pelo aluno.
\section{Referencial teórico}

O transporte aéreo apresenta características que tornam a gestão desse tipo de despesa particularmente desafiadora. Diferentemente de itens de aquisição com preços mais estáveis ou previamente contratados, as tarifas aéreas variam de forma dinâmica em função de múltiplos fatores, como antecedência da compra, origem e destino, sazonalidade, disponibilidade de assentos, concorrência entre companhias, dia da semana, horário do voo, eventos locais, nível de demanda e condições gerais do mercado. Dessa forma, o valor de uma passagem pode oscilar significativamente entre rotas, períodos e perfis de deslocamento, dificultando a adoção de critérios únicos e estáticos para orientar o planejamento das compras.

Além disso, a análise de passagens efetivamente compradas permite compreender padrões históricos de gasto e comportamento institucional, mas não revela, por si só, qual teria sido o preço da mesma passagem caso a aquisição tivesse ocorrido em outro momento. Essa limitação é relevante porque a avaliação da antecedência de compra envolve uma dimensão contrafactual: para afirmar que uma passagem teria sido mais barata em outro prazo, seria necessário observar ou estimar preços alternativos para condições comparáveis de origem, destino, data, horário, companhia aérea e características do voo. Assim, o uso exclusivo de registros de compras realizadas pode ser insuficiente para sustentar conclusões robustas sobre a variação de preços ao longo do tempo.

Para avaliar previsões de preço sem usar informação futura na aprendizagem, \citeonline[seção 5.10]{hyndman2021} apresentam validação temporal com conjuntos de treino anteriores a cada observação testada. Essa referência sustenta a organização temporal da avaliação proposta no pré-projeto, mas não estabelece o desempenho do modelo nas rotas desta pesquisa. \citeonline{breiman2001} apresenta Random Forests como combinação de preditores de árvores e explicita sua aplicação à regressão; isso fundamenta incluir essa família como alternativa para o preço, sem presumir superioridade sobre referências simples nem economia comprovada.

O pré-projeto define a comparação com regras simples como condição para considerar útil uma abordagem preditiva de antecedência. A seção 4.5 fundamenta essa decisão, em vez de escolher Random Forest apenas por sua sofisticação:
\begin{citacao}
Inicialmente, serão construídos modelos de referência, ou baselines, para permitir comparação com abordagens mais sofisticadas. Esses baselines poderão incluir média histórica por rota, mediana por faixa de antecedência, último preço observado, regra fixa de compra por antecedência ou recomendações baseadas em janelas administrativas comuns, como 10, 15, 21, 30 ou 45 dias. A comparação com baselines será fundamental para avaliar se o uso de técnicas preditivas efetivamente agrega valor ao processo decisório. \parencite[seção 4.5]{goncalves2026}
\end{citacao}

Assim, a inclusão de Random Forest é uma alternativa de comparação, e o valor agregado sobre uma regra simples permanece uma hipótese a testar nas cotações controladas. Não há nesta dissertação dummy resultado de economia comprovada.
